\begin{frame}{cf.\ 시계열이면 무작위로 쪼개면 안 된다}
  \begin{columns}[T]
    \begin{column}{0.42\textwidth}
      \centering
      \includegraphics[height=5.0cm]{ref_nested_cv_timeseries.png}
    \end{column}
    \begin{column}{0.56\textwidth}
      \footnotesize
      \begin{itemize}
        \item 일반 $k$-fold는 \textbf{무작위로 섞어서} 쪼갠다 ---
          시계열에서는 \textbf{미래 데이터가 과거 학습에 새는}
          것(data leakage)과 같다.
        \item \kb{Rolling-origin}(확장 윈도우): 항상 \textbf{앞
          구간만 train}, 그 다음 구간을 test로 --- 시간이 지날수록
          train 구간이 넓어지며 반복.
        \item 그림의 \kb{Outer loop}는 이렇게 test 오차를 여러 번
          측정해 평균내고, \kb{Inner loop}는 그 안에서(train
          구간만으로) 하이퍼파라미터를 튜닝 --- \textbf{중첩(nested)
          CV}.
      \end{itemize}
    \end{column}
  \end{columns}
\end{frame}
